\section{An Update in Place and a Sum of Remainders}\label{sec:conda}

\subsection{The program}
The test \texttt{sum/conda} fills an array with ones, modifies it in place by cases on the old value, and
sums it; the assertion is the exact sum.
\begin{lstlisting}[language=C]
sum[0] = 0;
for (i = 0; i < N; i++) a[i] = 1;
for (i = 0; i < N; i++) { if (a[i] == 1) a[i] = a[i] + 1; else a[i] = a[i] - 1; }
for (i = 0; i < N; i++) sum[0] = sum[0] + a[i];
__VERIFIER_assert(sum[0] == 2 * N);
\end{lstlisting}

\subsection{The trace equation system}
Three loops in sequence, each over its own counter of $n$, $n'$, $n''$ steps, all equal to $N$ by their
tests and exits as in Section~\ref{sec:brs2}. The middle one is the interesting equation: its update reads
the cell it writes, and its counter trace $I'$ has the array among its dependencies, since the guards
mention it,
\begin{align}
  A &= \tr{a_0}{\bl{K}, \bl{I}}{a, k, i \mapsto a[i := 1]} \label{eq:cd1}\\
  A' &= \tr{a'}{\bl{K'}, \bl{I'}}{a, k, i \mapsto \upd{a[i] = 1}{a[i := a[i] + 1]}\ \ \upd{a[i] \ne 1}{a[i := a[i] - 1]}} \label{eq:cd2}\\
  I' &= \tr{0}{\bl{K'}, \bl{A'}}{i, k, a \mapsto \upd{a[i] = 1}{i + 1}\ \ \upd{a[i] \ne 1}{i + 1}} \nonumber\\
  S &= \tr{s_0[0 := 0]}{\bl{K''}, \bl{I''}}{s, k, i \mapsto s[0 := s[0] + a''[i]]} \label{eq:cd3}\\
  &\Longrightarrow\ \cur(S)[0] = 2N \nonumber
\end{align}
with $a' = \cur(A)$ and $a'' = \cur(A')$.

\subsection{The lemmas}
\begin{align*}
  &\lemma{stores\_cases\_def}:\\
  &\quad \mathsf{trace}(a_0, \mathsf{for}\ \bl{\cnt(0, 1, n)}, \bl{\cnt(0, 1, n)}\ \mathsf{do}\ a, k, i \mapsto {}\\
  &\qquad\qquad \upd{P(a[k])}{a[k := c]}\ \ \upd{\neg P(a[k])}{a[k := G(a[k])]})\\
  &\quad = \mathsf{stores}(a_0, \bl{\cnt(0, 1, n)}, \mapc{\bl{\cnt(0, 1, n)}}{k \mapsto \ite{P(a_0[k])}{c}{G(a_0[k])}})
\end{align*}
Each position $k$ is written once, in iteration $k$, and read only there before it is written, so the loop
stores values computed from the \emph{initial} array. The first case stores a constant and not $G'(a[k])$: a
higher-order instance must contain the term it abstracts (\texttt{Rewrite.abstraction}), and $a[k] + 1$ is
$2$ in this program by the time the lemma is tried (step~2 below). The rest is the database's:
\lemma{stores\_def} and \lemma{maps\_const} for the first loop, \lemma{current\_stores} with
\lemma{occurs\_counter\_up}, \lemma{indexof\_counter\_up} and \lemma{earlier\_consts} for the reads,
\lemma{current\_trace\_add\_at} for the sum at the constant key, \lemma{selects\_stores\_counter} and
\lemma{sum\_consts} for the result.

\subsection{The proof}
Of the 69 rule applications:
\begin{enumerate}[nosep]
\item \textbf{The first loop.} \eqref{eq:cd1} is $\mathsf{stores}(a_0, \bl{I}, \mapc{\bl{K}}{k \mapsto 1})$ by
  \lemma{stores\_def}, and the map of a constant is $\bl{\mathsf{consts}(n, 1)}$ (\lemma{maps\_const}); $I$
  is $\cnt(0, 1, n)$.
\item \textbf{The guards of the second loop.} In the context of the iteration $k$ of \eqref{eq:cd2}, the
  previous value $a$ at position $k$ is unmodified, $a[k] = a'[k]$, and $a'[k]$ is read from the stores of the
  first loop: \lemma{current\_stores} gives
  $\ite{\mathsf{occurs}(k, \bl{I})}{\mathsf{earlier}(\mathsf{indexof}(k, \bl{I}), \bl{\mathsf{consts}(n, 1)})}{a_0[k]}$,
  the test is $0 \le k < \intf(n)$ by \lemma{occurs\_counter\_up}, true in the iteration, and the read of the
  constant sequence is $1$ (\lemma{earlier\_consts}). So $a[k] = 1$ inside the update: the first result
  $a[k] + 1$ is $2$, the guards are $a[k] = 1$ and $a[k] \ne 1$ on the \emph{bound} variable $a$, which the
  lemma's $P$ abstracts.
\item \textbf{The second loop is stores.} \lemma{stores\_cases\_def} matches \eqref{eq:cd2} with
  $P := \lambda z.\ z = 1$, $c := 2$, $G := \lambda z.\ z - 1$ and gives
  $\mathsf{stores}(a', \bl{\cnt(0, 1, n')}, \mapc{\bl{\cnt(0, 1, n')}}{k \mapsto \ite{a'[k] = 1}{2}{a'[k] - 1}})$;
  the read $a'[k]$ in the map is again $1$ by step~2, the conditional is $2$, and the map is
  $\bl{\mathsf{consts}(n', 2)}$.
\item \textbf{The sum.} \eqref{eq:cd3} is a trace that adds to the cell $0$ in every iteration;
  \lemma{current\_trace\_add\_at} makes $\cur(S) = s_0[0 := 0][0 := 0 + \mathsf{sum}(\mapc{\bl{K''}, \bl{I''}}{k, i \mapsto a''[i]})]$,
  the map of reads is $\mathsf{selects}(a'', \bl{I''})$ (\lemma{selects}), and with $n'' = n' = n$
  (\lemma{integer\_injective} from the three exits) \lemma{selects\_stores\_counter} reads the stored values
  back at the positions of the counter: $\bl{\mathsf{consts}(n, 2)}$, whose sum is $2\,\intf(n)$
  (\lemma{sum\_consts}). The postcondition $2\,\intf(n) = 2N$ holds by $\intf(n) = N$.
\end{enumerate}
The tests \texttt{s1lif} to \texttt{s5lif}, whose second loops have the same shape with other cases, close by
the same steps (69 rule applications for \texttt{s1lif}); a false assertion \texttt{sum[0] == 2 * N + 1}
leaves $2\,\intf(n) = 2N + 1$, i.e.\ the lemmas compute the sum rather than close the goal.

\subsection{A sum of remainders}\label{sec:ms2}
The test \texttt{sum/ms2} stores $i \bmod 2$ at $i$ and asserts \texttt{sum[0] <= N} for the sum of
Section~\ref{sec:brs2}. The first loop is plain stores of a map, $\mathsf{stores}(a_0, \bl{I}, \mapc{\bl{K}, \bl{I}}{k, i \mapsto \ite{0 \le i}{i \bmod 2}{-(-i \bmod 2)}})$
(the remainder of C under a sign that the context of the map does not know yet), the read $a'[k]$ in the
second loop goes through \lemma{current\_stores} to the map at $k$, where $0 \le k$ holds and the read is
$k \bmod 2$, \lemma{select\_current\_trace\_reset\_add} abstracts $G := \lambda k.\ k \bmod 2$, and
\begin{align*}
  &\lemma{sum\_mod\_below}:\quad 0 < m\ \wedge\ \intf(n) \cdot (m - 1) \le e\\
  &\quad\Longrightarrow\ \mathsf{sum}(\mapc{\bl{\cnt(0, 1, n)}}{k \mapsto \ite{0 < k}{k \bmod m}{0}}) \le e\ \Longleftrightarrow\ \mathsf{true}
\end{align*}
closes the bound with $m := 2$, $e := N$ and $\intf(n) = N$ (61 rule applications). The three tests
\texttt{brs2}, \texttt{conda} and \texttt{ms2} share every step but the last two: the shape of the summands
(constants by index, constants by value, remainders) selects the introduction lemma for the stores and the
bound lemma for the sum.

\subsection{The idea}
Loops in sequence compose through the current values, as in the counting sort (Section~\ref{sec:cs}): the
stores of one loop are read in the update of the next, where the context of the iteration decides the
membership tests, and the read collapses to the stored value before the introduction lemma of the next loop
is tried. The order matters: \lemma{stores\_cases\_def} sees $a[k] + 1$ as $2$ because the read was resolved
first, which is why its first case is a constant; a lemma written from the program text, with $G'(a[k])$ in
the first case, would not match the simplified update. The lemmas for a program are written for the goal as
the simplifier prints it, not for the program (the pattern of Section~\ref{sec:principles:matcher}, arithmetic
in the goal's form).
